% GENERATED FILE. Edit canonical lessons, then run npm run generate:books.
% canonical-source-hash: fnv1a64:658d14d81f296a81
% canonical-lessons: SA-C544-parasparam, SA-C544-yugapat, SA-C544-mithah, SA-C544-prthak, SA-C544-samprati

\chapter{Each in Turn}
\label{ch:sa-each-in-turn}
\hlchaptermodality{544}

\begin{chapteropening}
\textbf{By the end of this chapter:} \emph{I can say whether people act together, apart or with each other.}

Complete the last lesson of chapter 544: I can say whether people act together, apart or with each other.
\end{chapteropening}

\section[\texorpdfstring{\sk{परस्परम्} (parasparam)}{परस्परम् (parasparam)}]{\sk{परस्परम्} (parasparam) — each other}

\label{lesson:SA-C544-parasparam}



\begin{quote}
Before the new one: say the Sanskrit for insistence, then the Sanskrit for by oneself.
\end{quote}

\subsection*{You'll want to know: \sk{परस्परम्}}
\textbf{\sk{परस्परम्}} — \emph{parasparam} — \textquotedblleft{}each other\textquotedblright{}.

\begin{quote}
\sk{यदि} \sk{भ्रातरौ} \sk{परस्परं} \sk{स्निह्यतः}, \sk{तर्हि} \sk{माता} \sk{प्रसन्ना} \sk{भवति।}
\end{quote}

\begin{quote}
\emph{yadi bhrātarau parasparaṃ snihyataḥ, tarhi mātā prasannā bhavati.}
\end{quote}

\begin{quote}
If the two brothers love each other, my mother is happy.
\end{quote}

\begin{cousinweb}
\textbf{\sk{परस्परम्}} (\emph{parasparam}) is built from \emph{para}, \textquotedblleft{}other\textquotedblright{}, said twice, the first time in its nominative \emph{paras}; MW gives \emph{paraspara} as mutual, each other's (MW).
\end{cousinweb}

\begin{grammarlens}[title={What you've built}]
One more word for saying what happens if.
\end{grammarlens}

\subsection*{Your turn}
\begin{itemize}
\raggedright
  \item \emph{Say it:} \emph{parasparam}
  \item \emph{Say it:} \emph{parasparam}, once more
  \item \emph{Say it:} \emph{yadi bhrātarau parasparaṃ snihyataḥ, tarhi mātā prasannā bhavati.} — If the two brothers love each other, my mother is happy.
  \item \emph{Recall:} say the Sanskrit for insistence, then the Sanskrit for by oneself, then say \emph{parasparam} again
\end{itemize}

\begin{tcolorbox}[breakable,colback=teal!4,colframe=teal!35!black,title={Before you move on}]
What does \sk{परस्परम्} mean? (\textquotedblleft{}Each other\textquotedblright{}.) Say it once more.
\end{tcolorbox}
\section[\texorpdfstring{\sk{युगपत्} (yugapat)}{युगपत् (yugapat)}]{\sk{युगपत्} (yugapat) — at the same time}

\label{lesson:SA-C544-yugapat}



\begin{quote}
Before the new one: say the Sanskrit for by oneself, then the Sanskrit for each other.
\end{quote}

\subsection*{You'll want to know: \sk{युगपत्}}
\textbf{\sk{युगपत्}} — \emph{yugapat} — \textquotedblleft{}at the same time\textquotedblright{}.

\begin{quote}
\sk{यदि} \sk{सर्वे} \sk{युगपत्} \sk{वदन्ति}, \sk{तर्हि} \sk{अहं} \sk{न} \sk{अवगच्छामि।}
\end{quote}

\begin{quote}
\emph{yadi sarve yugapat vadanti, tarhi ahaṃ na avagacchāmi.}
\end{quote}

\begin{quote}
If everyone talks at the same time, I don't understand.
\end{quote}

\begin{cousinweb}
\textbf{\sk{युगपत्}} (\emph{yugapat}) is built from \emph{yuga}, \textquotedblleft{}a yoke\textquotedblright{}, and \emph{pad}, \textquotedblleft{}a foot\textquotedblright{}: side by side in one yoke; MW gives \emph{yugapad} as together, at the same time (MW).
\end{cousinweb}

\begin{grammarlens}[title={What you've built}]
One more word for saying what happens if.
\end{grammarlens}

\subsection*{Your turn}
\begin{itemize}
\raggedright
  \item \emph{Say it:} \emph{yugapat}
  \item \emph{Say it:} \emph{yugapat}, once more
  \item \emph{Say it:} \emph{yadi sarve yugapat vadanti, tarhi ahaṃ na avagacchāmi.} — If everyone talks at the same time, I don't understand.
  \item \emph{Recall:} say the Sanskrit for by oneself, then the Sanskrit for each other, then say \emph{yugapat} again
\end{itemize}

\begin{tcolorbox}[breakable,colback=teal!4,colframe=teal!35!black,title={Before you move on}]
What does \sk{युगपत्} mean? (\textquotedblleft{}At the same time\textquotedblright{}.) Say it once more.
\end{tcolorbox}
\section[\texorpdfstring{\sk{मिथः} (mithaḥ)}{मिथः (mithaḥ)}]{\sk{मिथः} (mithaḥ) — together; mutually}

\label{lesson:SA-C544-mithah}



\begin{quote}
Before the new one: say the Sanskrit for each other, then the Sanskrit for at the same time.
\end{quote}

\subsection*{You'll want to know: \sk{मिथः}}
\textbf{\sk{मिथः}} — \emph{mithaḥ} — \textquotedblleft{}together; mutually\textquotedblright{}.

\begin{quote}
\sk{यदि} \sk{बालाः} \sk{मिथः} \sk{क्रीडन्ति}, \sk{तर्हि} \sk{ते} \sk{प्रसन्नाः} \sk{सन्ति।}
\end{quote}

\begin{quote}
\emph{yadi bālāḥ mithaḥ krīḍanti, tarhi te prasannāḥ santi.}
\end{quote}

\begin{quote}
If the children play together, they are happy.
\end{quote}

\begin{cousinweb}
\textbf{\sk{मिथः}} (\emph{mithaḥ}) is the word MW gives for together, mutually, alternately; MW gives no derivation, so its origin is uncertain (MW).
\end{cousinweb}

\begin{grammarlens}[title={What you've built}]
One more word for saying what happens if.
\end{grammarlens}

\subsection*{Your turn}
\begin{itemize}
\raggedright
  \item \emph{Say it:} \emph{mithaḥ}
  \item \emph{Say it:} \emph{mithaḥ}, once more
  \item \emph{Say it:} \emph{yadi bālāḥ mithaḥ krīḍanti, tarhi te prasannāḥ santi.} — If the children play together, they are happy.
  \item \emph{Recall:} say the Sanskrit for each other, then the Sanskrit for at the same time, then say \emph{mithaḥ} again
\end{itemize}

\begin{tcolorbox}[breakable,colback=teal!4,colframe=teal!35!black,title={Before you move on}]
What does \sk{मिथः} mean? (\textquotedblleft{}Together; mutually\textquotedblright{}.) Say it once more.
\end{tcolorbox}
\section[\texorpdfstring{\sk{पृथक्} (pṛthak)}{पृथक् (pṛthak)}]{\sk{पृथक्} (pṛthak) — separately}

\label{lesson:SA-C544-prthak}



\begin{quote}
Before the new one: say the Sanskrit for at the same time, then the Sanskrit for together.
\end{quote}

\subsection*{You'll want to know: \sk{पृथक्}}
\textbf{\sk{पृथक्}} — \emph{pṛthak} — \textquotedblleft{}separately\textquotedblright{}.

\begin{quote}
\sk{यदि} \sk{पुस्तकानि} \sk{पृथक्} \sk{सन्ति}, \sk{तर्हि} \sk{तानि} \sk{पेटिकायां} \sk{स्थापयतु।}
\end{quote}

\begin{quote}
\emph{yadi pustakāni pṛthak santi, tarhi tāni peṭikāyāṃ sthāpayatu.}
\end{quote}

\begin{quote}
If the books are lying about separately, put them in the box.
\end{quote}

\begin{cousinweb}
\textbf{\sk{पृथक्}} (\emph{pṛthak}) is built on the root √pṛth (or \emph{prath}), \textquotedblleft{}to spread\textquotedblright{}, and \emph{añc}, \textquotedblleft{}to turn\textquotedblright{}, as MW gives it: spread apart; MW gives \emph{pṛthak} as widely apart, separately, singly (MW).
\end{cousinweb}

\begin{grammarlens}[title={What you've built}]
One more word for saying what happens if.
\end{grammarlens}

\subsection*{Your turn}
\begin{itemize}
\raggedright
  \item \emph{Say it:} \emph{pṛthak}
  \item \emph{Say it:} \emph{pṛthak}, once more
  \item \emph{Say it:} \emph{yadi pustakāni pṛthak santi, tarhi tāni peṭikāyāṃ sthāpayatu.} — If the books are lying about separately, put them in the box.
  \item \emph{Recall:} say the Sanskrit for at the same time, then the Sanskrit for together, then say \emph{pṛthak} again
\end{itemize}

\begin{tcolorbox}[breakable,colback=teal!4,colframe=teal!35!black,title={Before you move on}]
What does \sk{पृथक्} mean? (\textquotedblleft{}Separately\textquotedblright{}.) Say it once more.
\end{tcolorbox}
\section[\texorpdfstring{\sk{सम्प्रति} (samprati)}{सम्प्रति (samprati)}]{\sk{सम्प्रति} (samprati) — at present}

\label{lesson:SA-C544-samprati}



\begin{quote}
Before the new one: say the Sanskrit for together, then the Sanskrit for separately.
\end{quote}

\subsection*{You'll want to know: \sk{सम्प्रति}}
\textbf{\sk{सम्प्रति}} — \emph{samprati} — \textquotedblleft{}at present\textquotedblright{}.

\begin{quote}
\sk{यदि} \sk{सम्प्रति} \sk{वर्षा} \sk{भवति}, \sk{तर्हि} \sk{वयं} \sk{पश्चात्} \sk{गमिष्यामः।}
\end{quote}

\begin{quote}
\emph{yadi samprati varṣā bhavati, tarhi vayaṃ paścāt gamiṣyāmaḥ.}
\end{quote}

\begin{quote}
If it is raining now, we will go later.
\end{quote}

\begin{cousinweb}
\textbf{\sk{सम्प्रति}} (\emph{samprati}) is built on \emph{sam} and \emph{prati}, \textquotedblleft{}towards, against\textquotedblright{}; MW gives \emph{samprati} as directly opposite, at the right time, and now, at present (MW).
\end{cousinweb}

\begin{grammarlens}[title={What you've built}]
One more word for saying what happens if.
\end{grammarlens}

\subsection*{Your turn}
\begin{itemize}
\raggedright
  \item \emph{Say it:} \emph{samprati}
  \item \emph{Say it:} \emph{samprati}, once more
  \item \emph{Say it:} \emph{yadi samprati varṣā bhavati, tarhi vayaṃ paścāt gamiṣyāmaḥ.} — If it is raining now, we will go later.
  \item \emph{Recall:} say the Sanskrit for together, then the Sanskrit for separately, then say \emph{samprati} again
\end{itemize}

\begin{tcolorbox}[breakable,colback=teal!4,colframe=teal!35!black,title={Before you move on}]
What does \sk{सम्प्रति} mean? (\textquotedblleft{}At present\textquotedblright{}.) Say it once more.
\end{tcolorbox}
